\documentclass[11pt,a4paper]{article}
\usepackage{amsmath,amssymb,bm}
\usepackage[margin=2.2cm]{geometry}
\usepackage{parskip}
\usepackage{microtype}
\usepackage{booktabs}
\usepackage{listings}

\setlength{\parindent}{0pt}
\lstset{basicstyle=\ttfamily\small, frame=single, aboveskip=4pt, belowskip=4pt}

\title{\vspace{-1.2cm}\textbf{AVBD's Explosive-Error-Correction Guard Is Disabled\\ for Joints Touching an Elastic Endpoint}
\\[0.3em]\large\normalfont Why a clamped reduced elastic body diverges at small steps and few iterations}
\author{\normalsize Newton Physics Engine}
\date{\normalsize August 2026}

\begin{document}
\maketitle
\vspace{-0.6cm}

\section{The observation}

A reduced elastic blade clamped to a heavier carriage diverges in free fall, with no contact,
when the substep count is high and the iteration count low. The stable region is bounded from
\emph{below} in iterations, and the bound rises with the substep count:

\begin{center}\small
\begin{tabular}{@{}rccccc@{}}
\toprule
substeps & 2 iter & 5 iter & 10 iter & 20 iter & 40 iter \\
\midrule
4  & ok  & ok  & ok  & ok  & ok \\
8  & div & ok  & ok  & ok  & ok \\
12 & div & div & ok  & ok  & ok \\
16 & div & div & ok  & div & ok \\
24 & div & div & div & ok  & ok \\
32 & div & div & div & div & div \\
\bottomrule
\end{tabular}
\end{center}

Refining the step normally reduces the iterations needed, which is the argument of Macklin et
al., \emph{Small Steps in Physics Simulation} (2019). Here it raises them.

\section{The defect}

AVBD regularises the constraint function itself. Giles, Diaz and Yuksel,
\emph{Augmented Vertex Block Descent}, SIGGRAPH 2025, Section 3.6, define
\begin{equation}
  C_j(x) = C^*_j(x) - \alpha\,C^*_j(x^t),
  \qquad \alpha = 0.95 ,
  \label{eq:reg}
\end{equation}
with $C^*_j$ the raw constraint and $x^t$ the configuration at the start of the timestep. The
parameter $\alpha$ sets what portion of the pre-existing constraint error is ignored. The
section is titled \emph{Preventing Explosive Error Correction}, and its motivation is stated
directly:

\begin{quote}\itshape
if a hard constraint is violated in the previous frame (due to limited iterations), in the very
next frame the hard constraint can apply an arbitrarily large force to enforce the constraint,
resulting in a sudden motion and a large spike in momentum.
\end{quote}

The regularised $C_j$ enters the energy, the force and the dual update alike. With the dual
update of the paper's Equation 11, $\lambda^{(n+1)} = k^{(n)} C_j(x) + \lambda^{(n)}$, the
implementation is (\texttt{rigid\_vbd\_kernels.py:353})
\begin{equation}
  \bm{\lambda}_{n+1} = \bm{\lambda}_n + k\,\big(\mathbf{C}_n - \alpha\,\mathbf{C}_0\big),
  \qquad
  \mathbf{f} = k\,\mathbf{C} + \bm{\lambda},
  \label{eq:dual}
\end{equation}
where $\mathbf{C}_0$ is the snapshot of $C^*_j(x^t)$. This is \eqref{eq:reg}, with the paper's
value of $\alpha$.

The snapshot kernel skips exactly the joints of interest
(\texttt{step\_joint\_C0\_lambda:2145}):

\begin{lstlisting}
if joint_parent_elastic_endpoint[j] >= 0 or joint_child_elastic_endpoint[j] >= 0:
    joint_C0_lin[j] = wp.vec3(0.0)
    joint_C0_ang[j] = wp.vec3(0.0)
    joint_lambda_lin[j] = joint_lambda_lin[j] * lambda_decay
    joint_lambda_ang[j] = joint_lambda_ang[j] * lambda_decay
    return
\end{lstlisting}

Any joint touching an elastic endpoint gets $\mathbf{C}_0 = \mathbf{0}$ and returns before the
snapshot below. Setting $\mathbf{C}_0 = \mathbf{0}$ is setting $\alpha = 0$ in \eqref{eq:reg}:
no pre-existing error is ignored, and the joint is placed in exactly the configuration
Section~3.6 of the paper exists to prevent. The dual is left active and unregularised,
$\bm{\lambda} \mathrel{+}= k\,\mathbf{C}$, integrating the full standing violation every
iteration.

One part of the paper's scheme does survive the branch: $\bm{\lambda}$ is still scaled down
between steps by $\alpha\gamma$, the warm-start factor of the paper's Equation 19, visible as
\texttt{lambda\_decay} in the listing and computed as
\texttt{rigid\_joint\_alpha * rigid\_avbd\_gamma} in the solver. That decay is then the only
thing left bounding the multiplier. Section~\ref{sec:alpha0} shows it is enough when the driving
term is the small residue that \eqref{eq:reg} leaves behind, and not enough once
\eqref{eq:reg} is removed and the multiplier is driven by the full violation.

Measured over 256 substeps on the clamp joint, comparing the elastic model against a rigid pair
of the same masses and topology. The rigid pair is driven by a $500$ N force on the light body
oscillating at $8$ Hz, so that its clamp carries a time-varying load and its multiplier is
genuinely engaged; without such a drive the violation is static, the snapshot cancels it exactly
and $\bm{\lambda}$ never leaves zero, which makes the comparison vacuous.

\begin{center}\small
\begin{tabular}{@{}lrrr@{}}
\toprule
 & rms $\lvert\mathbf{C}\rvert$ & rms $\lvert\mathbf{C}_0\rvert$ & rms $\lvert\mathbf{C}-\alpha\mathbf{C}_0\rvert$ \\
\midrule
elastic clamp & $7.82\times10^{-8}$ & $0$ exactly      & $7.82\times10^{-8}$ \\
rigid clamp   & $1.733\times10^{-4}$ & $1.724\times10^{-4}$ & $9.49\times10^{-6}$ \\
\bottomrule
\end{tabular}
\end{center}

For the rigid clamp the snapshot tracks the violation and the stabilisation removes $95\%$ of
it. For the elastic clamp nothing is removed.

\subsection{Why textbook ALM does not need this}

The augmented Lagrangian method carries no such term and converges unconditionally, so the
regularisation can look like an unnecessary graft. The guarantee, however, is for a single
static problem whose primal is minimised before each dual update. AVBD runs the method inside a
timestepping loop, with a handful of iterations per step and $\lambda$ warm-started across
steps. Sections 3.6 and 3.7 of the paper exist because that regime injects energy: pre-existing
error drives an arbitrarily large correction, and warm-starting carries it forward. The
$\alpha$ regularisation is not a departure from ALM but the correction that makes ALM usable
here, and the paper says as much.

\section{The rigid control}

Forcing $\mathbf{C}_0 = \mathbf{0}$ on the rigid pair, changing nothing else, reproduces the
failure --- with no elastic body, no modal coordinate and no floating frame anywhere in the
model:

\begin{center}\small
\begin{tabular}{@{}rrrll@{}}
\toprule
substeps & iterations & $k$ & normal $\mathbf{C}_0$ & $\mathbf{C}_0$ zeroed \\
\midrule
24 & 2  & $10^{6}$ & ok, $\lvert\bm{\lambda}\rvert \to 538$ & diverged \\
24 & 20 & $10^{6}$ & ok, $\lvert\bm{\lambda}\rvert \to 500$ & diverged \\
24 & 2  & $10^{7}$ & ok, $\lvert\bm{\lambda}\rvert \to 520$ & diverged \\
4  & 2  & $10^{6}$ & ok, $\lvert\bm{\lambda}\rvert \to 704$ & diverged \\
\bottomrule
\end{tabular}
\end{center}

Zeroing the snapshot is sufficient on its own. In this pure form the failure is also
unconditional: neither more iterations, nor a stiffer penalty, nor a coarser step recovers it.

\section{What $\alpha = 0$ does to the multiplier}
\label{sec:alpha0}

Every term of \eqref{eq:dual} was recorded per iteration on the rigid pair. With $N$ iterations
per step and the decay measured at $d = 0.949$, one step is
\begin{equation}
  \lambda_{n+1} = d\Big(\lambda_n + k\sum_{i=1}^{N} C_i\Big),
  \label{eq:step}
\end{equation}
and \eqref{eq:dual} reproduces the recorded multiplier to a relative residual of $10^{-7}$.

\subsection{Proportional against integral feedback}

The snapshot is the violation left at the end of the previous step: measured over $179$ steps,
$\mathbf{C}_0$ of step $n$ equals $C$ at the end of step $n-1$ to a median relative difference of
$2\times10^{-5}$. Writing $C_n$ for that violation and taking one update per step for clarity,
\eqref{eq:dual} reads
\begin{equation}
  \lambda_{n+1} = \lambda_n + k\,\big(C_n - \alpha\,C_{n-1}\big).
  \label{eq:diff}
\end{equation}

At $\alpha \to 1$ the right-hand side is a difference and the sum telescopes to
$\lambda_{n+1} = \lambda_0 + k\,C_n - k\,C_{-1}$: the multiplier \emph{tracks} the violation,
bounded by $kC$ however long the simulation runs. At $\alpha = 0$ the sum no longer telescopes,
\begin{equation}
  \lambda_{n+1} = \lambda_n + k\,C_n = \lambda_0 + k\sum_{i=0}^{n} C_i ,
\end{equation}
and the multiplier \emph{integrates} the violation instead. Proportional feedback becomes
integral feedback.

This is the paper's explosive error correction seen from the other side: a multiplier that
integrates the standing violation is one that keeps adding force to correct an error the step
began with. It also predicts the form of the failure. An integrator contributes ninety degrees
of phase lag, so the loop should fail by growing oscillation rather than by monotone growth.
Section~\ref{sec:mode} shows that it does.

At the value actually used, $\alpha = 0.95$, \eqref{eq:diff} is a leaky difference: $95\%$
telescopes and a $5\%$ residue accumulates against the decay. That residue sets a finite fixed
point. With $C_n \approx C_{n-1} \equiv C$, each update contributes $k(1-\alpha)C$, so
\begin{equation}
  \lambda_{n+1} = d\big(\lambda_n + N k (1-\alpha) C\big),
  \qquad
  \lambda^\star = \frac{d\,N\,k\,(1-\alpha)\,C}{1-d} .
  \label{eq:fixed}
\end{equation}
With $d = 0.949$, $N = 2$, $k = 10^{6}$, $\alpha = 0.95$ and the measured
$C = 2.36\times10^{-4}$, this gives $22.4/0.051 \approx 439$ against a recorded $\lambda$
oscillating between $466$ and $491$: within seven percent. The quasi-static assumption behind
\eqref{eq:fixed} holds only in the stable regime; once the mode below is growing, $C$ oscillates
at nine steps per period and $C_n \not\approx C_{n-1}$.

\subsection{The mode this produces}
\label{sec:mode}

Everything measured here is a property of the $\lambda$ time series.

$\lambda$ fails by oscillation rather than by growth. Under the sustained drive it runs
$+428 \to -75651 \to +1.5\times10^{7}$, changing sign, which no scalar real recurrence can
produce.

Removing the drive isolates the mode. Applying a single six-step impulse and then holding the
force at zero for the rest of the run, the free evolution of $\lambda$ is:

\begin{center}\small
\begin{tabular}{@{}lrrrr@{}}
\toprule
 & sign changes & period & frequency & AR(2) residual \\
\midrule
$\alpha = 0.95$ & 27 & $66.5$ steps & $72.2$ Hz & $1.1\times10^{-1}$ \\
$\alpha = 0$    & 54 & $8.6$ steps  & $558.8$ Hz & $8.0\times10^{-7}$ \\
\bottomrule
\end{tabular}
\end{center}

Both cases ring, but there is no forcing left to ring at: the frequency is a property of the
coupled multiplier--body system, and it differs by a factor of eight between the two. Only the
unregularised case grows, reaching $9.2\times10^{13}$ in $253$ steps from that single kick. An
AR(2) fitted to $\lambda$ reproduces its whole response to $8\times10^{-7}$, confirming a
second-order state, with the complex pair at
\begin{equation}
  \lvert\mu\rvert = 1.1100 , \qquad f = 558.8\text{ Hz} .
\end{equation}
Raising the model order leaves the dominant root unchanged. No second-order model fits the
regularised case, whose roots are a real $0.99$ and a pair at $0.856$, all inside the unit
circle --- consistent with a multiplier slaved to the violation rather than free to ring.

\subsection{The regularisation is causal}

The two estimates above are fits to one time series, so they establish the mode, not its cause.
Scaling only the clamp's snapshot, $\mathbf{C}_0 \mapsto \beta\,\mathbf{C}_0$, varies the
effective $\alpha$ continuously while leaving the decay untouched:

\begin{center}\small
\begin{tabular}{@{}rrrrl@{}}
\toprule
$\beta$ & effective $\alpha$ & $\lvert\mu\rvert$ & frequency & AR(2) residual \\
\midrule
$1.0$ & $0.950$ & $0.8969$ & $177$ Hz & $1.1\times10^{-1}$ \\
$0.7$ & $0.665$ & $0.9841$ & $438$ Hz & $6.6\times10^{-3}$ \\
$0.5$ & $0.475$ & $1.0323$ & $476$ Hz & $6.1\times10^{-8}$ \\
$0.0$ & $0.000$ & $1.1100$ & $559$ Hz & $5.8\times10^{-8}$ \\
\bottomrule
\end{tabular}
\end{center}

The modulus rises monotonically as the regularisation is weakened, crosses unity near
$\beta \approx 0.6$, and terminates at the value measured with the snapshot zeroed.

Sweeping $\alpha$ itself instead would report stability throughout and isolate nothing:
\texttt{solver\_vbd.py:2624} sets
\texttt{joint\_lambda\_decay = rigid\_joint\_alpha * rigid\_avbd\_gamma}, the paper's Equation
19, so lowering $\alpha$ also drives the warm-start factor toward zero, wiping $\lambda$ at the
start of every step and preventing the multiplier accumulating at all.

\subsection{Where the boundary falls}

Substeps, iterations, penalty and mass ratio move the boundary without restoring the
regularisation; they change how large $C$ is and how quickly the primal retires it. Empirically
the boundary is set by $a = k h^2 / M$, the penalty against the inertial stiffness with $M$ the
reduced mass across the constraint. At 24 substeps and 2 iterations, varying $k$ gives
divergence at $a = 0.282$ and stability at $a = 0.326$; varying the mass ratio gives stability at
$a = 0.289$ and divergence at $a = 0.241$, intersecting at $a_{\text{crit}} \approx 0.28$.
Refining the step at fixed $k$ drives $a$ down, which is why the iteration requirement rises
rather than falls.

\section{Reproduction}

Three bodies, no contact: world $\to$ prismatic $\to$ carriage ($12.5$ kg) $\to$ fixed clamp
$\to$ reduced elastic blade ($0.3$ kg). Released from $5$ m, $0.35$ s, 24 substeps, 2
iterations.

The rigid control is the same topology with the blade replaced by a rigid body of equal mass,
driven by an oscillating force so the clamp is loaded, with \texttt{joint\_C0\_lin} and
\texttt{joint\_C0\_ang} zeroed after each snapshot.

\end{document}
